\clearpage
\section*{Сборка и воспроизведение}
\addcontentsline{toc}{section}{Сборка и воспроизведение}
Для C++ нужны компилятор с C++20, CMake 3.24+ и \code{just}.
Для Python используются \code{uv}, файлы \code{pyproject.toml}
и \code{uv.lock} каждого проекта. Для живого робота нужны Docker
и MongoDB replica set. Первая настройка тестов может требовать
сети для загрузки закреплённого Google Test.

Обычная сборка \code{just build} не строит индекс.
Команда \code{just build-index} - отдельная дорогостоящая
операция; для уже сохранённого индекса она не нужна.
Относительные пути в конфигурации считаются от каталога проекта.
Пример поиска:
\begin{lstlisting}
cd 13_tfidf
just build
printf 'memory allocation\n"memory allocation" / 5\n' |
  just search
just web
\end{lstlisting}
Команда тестирования:
\begin{lstlisting}
cd 17_unit_tests
just test
ctest --test-dir build --output-on-failure
\end{lstlisting}
Первая команда собирает и запускает тесты, вторая повторяет их
без пересборки. Для всех C++-проектов принцип тот же.

Измерение сжатия в лабораторной №11 выполняется после отдельного
построения VB- и raw-индексов командами \code{just build-index}
и \code{just build-raw}, затем \code{just measure}.
При уже существующем raw-индексе повторная подготовка завершается
ошибкой и не удаляет готовые файлы.

В лабораторной №12 команда \code{just prepare-tune} подготавливает
шесть вариантов индекса. Затем \code{just tune} измеряет их и
выбирает шаг. Для измерений корпус не нужен; отсутствие варианта
вызывает ошибку, а не автоматическую перестройку.
В лабораторной №13 \code{just measure} использует основной готовый
индекс и сравнивает строгий и ранжированный поиск.

Для пересчёта качества из каталога лабораторной №6:
\begin{lstlisting}
just evaluate-full
just plot-full
\end{lstlisting}
Первая команда проверяет согласованность запросов, выдач и оценок,
а затем вычисляет метрики. Вторая дополнительно строит графики.
Для этих команд не требуются корпус и поисковые индексы:
они работают по зафиксированным выдачам и разметке.
Неоценённый результат или изменение проверенных входных данных
приводит к диагностической ошибке.

\clearpage
\subsection*{Запросы и информационные потребности}
\begin{longtable}{lp{0.28\textwidth}p{0.51\textwidth}}
\toprule ID & Запрос & Информационная потребность\\\midrule\endfirsthead
\toprule ID & Запрос & Информационная потребность\\\midrule\endhead
q01 & \code{memory allocation} & Как выделять память в Linux и libc?\\
q02 & \code{memory allocations} & Тот же запрос во множественном числе; устойчивость к словоформе.\\
q03 & \code{mutex locking} & Как блокировать mutex в ядре и pthreads?\\
q04 & \code{spinlocks} & Как работают spinlock и когда их применять?\\
q05 & \code{workqueues} & Как планировать отложенную работу через workqueue?\\
q06 & \code{reference counting} & Как вести счётчик ссылок и не допустить use-after-free?\\
q07 & \code{memory mappings} & Как отображать и снимать отображение памяти?\\
q08 & \code{mprotect} & Как изменить права доступа к страницам памяти?\\
q09 & \code{fork} & Как создать дочерний процесс через fork?\\
q10 & \code{execve} & Как заменить образ процесса через execve?\\
q11 & \code{waiting children} & Как дождаться завершения дочерних процессов?\\
q12 & \code{pipes} & Как передавать данные через pipe?\\
q13 & \code{file descriptors} & Как открывать, дублировать и обслуживать файловые дескрипторы?\\
q14 & \code{synchronous writes} & Как гарантировать запись данных на носитель?\\
q15 & \code{epoll events} & Как получать уведомления готовности epoll?\\
q16 & \code{sockets} & Как создать и использовать сетевые сокеты?\\
q17 & \code{accept connections} & Как принимать входящие соединения?\\
q18 & \code{futex} & Как устроена синхронизация через futex?\\
q19 & \code{CPU affinity} & Как привязать задачу или прерывание к CPU?\\
q20 & \code{resource limits} & Как установить ограничения ресурсов процесса?\\
q21 & \code{pthread_create} & Как создать POSIX-поток? Точное имя API.\\
q22 & \code{malloc} & Как выделять память через malloc? Точное имя API.\\
q23 & \code{memory copies} & Как копировать блоки памяти?\\
q24 & \code{memory overlapping} & Как безопасно копировать перекрывающиеся области памяти?\\
q25 & \code{printf} & Как форматировать вывод? Точное имя API.\\
q26 & \code{dlopen} & Как загрузить разделяемую библиотеку? Точное имя API.\\
q27 & \code{huge pages} & Как использовать huge pages?\\
q28 & \code{cgroup memory} & Как ограничить память через cgroup v2?\\
q29 & \code{kernel sanitizers} & Как обнаружить ошибки памяти и гонки в ядре?\\
q30 & \code{epolll} & Опечатка в epoll; ожидается пустая выдача без коррекции.\\
\bottomrule
\end{longtable}

